\uuid{3fK8}
\titre{Développements en série entière au voisinage de $0$ }

\niveau{L2}
\module{Séries entières}
\chapitre{Séries entières}
\sousChapitre{Développement en série entière et rayon de convergence}
\theme{Calcul de développements en série entière en $0$}
\auteur{Quentin Liard}
\datecreate{ 2026-04-10}
\organisation{AMSCC}
\difficulte{2}

\contenu{

\texte{ 
On considère les fonctions
$$
f(x)=\frac{1}{3+\frac{x}{2}}
\qquad\text{et}\qquad
g(x)=\frac{\ln(1+x)}{x}.
$$

}

\begin{enumerate}
\item \question{Déterminer le développement en série entière en $0$ de la fonction $g$ puis donner son rayon de convergence.}
\indication{}
\reponse{
On sait que, pour $|x|<1$,
$$
\ln(1+x)=\sum_{n=1}^{+\infty}\frac{(-1)^{n+1}}{n}x^n.
$$
En divisant par $x$ (pour $x\neq 0$), on obtient
$$
\frac{\ln(1+x)}{x}
=
\sum_{n=1}^{+\infty}\frac{(-1)^{n+1}}{n}x^{n-1}.
$$
En posant $m=n-1$, cela devient
$$
\frac{\ln(1+x)}{x}
=
\sum_{m=0}^{+\infty}\frac{(-1)^m}{m+1}x^m.
$$
Donc
$$
\boxed{\frac{\ln(1+x)}{x}=\sum_{n=0}^{+\infty}\frac{(-1)^n}{n+1}x^n.}
$$

La série entière de $\ln(1+x)$ a pour rayon de convergence $1$, car sa singularité la plus proche de $0$ est en $x=-1$. La division par $x$ ne modifie pas le rayon de convergence de la série obtenue.

Ainsi,
$$
\boxed{R_g=1.}
$$
}
\item   \question{Déterminer le développement en série entière en $0$ de la fonction $h$ définie par
$$
h(x)=\frac{1}{1+x},
$$
puis donner son rayon de convergence.}
\indication{}
\reponse{
On reconnaît la série géométrique :
$$
\frac{1}{1-t}=\sum_{n=0}^{+\infty} t^n,
\qquad |t|<1.
$$
En prenant $t=-x$, on obtient
$$
\frac{1}{1+x}=\sum_{n=0}^{+\infty}(-1)^n x^n,
\qquad |x|<1.
$$
Donc
$$
\boxed{h(x)=\sum_{n=0}^{+\infty}(-1)^n x^n.}
$$

La singularité la plus proche de $0$ est en $x=-1$, donc le rayon de convergence est
$$
\boxed{R_h=1.}
$$
}
\item \question{En déduire le développement en série entière en $0$ de $f$ et également de la fonction $x \mapsto f(x)^2$.}
\reponse{
On commence par écrire
$$
f(x)=\frac{1}{3+\frac{x}{2}}
=\frac{1}{3\left(1+\frac{x}{6}\right)}
=\frac13\,h\left(\frac{x}{6}\right).
$$
En utilisant le développement en série entière de $h$, on obtient
$$
f(x)=\frac13\sum_{n=0}^{+\infty}(-1)^n\left(\frac{x}{6}\right)^n
=
\sum_{n=0}^{+\infty}\frac{(-1)^n}{3\cdot 6^n}x^n,
\qquad |x|<6.
$$
Ainsi
$$
\boxed{f(x)=\sum_{n=0}^{+\infty}\frac{(-1)^n}{3\cdot 6^n}x^n.}
$$

Pour $f(x)^2$, on a
$$
f(x)^2=\frac{1}{\left(3+\frac{x}{2}\right)^2}
=\frac{1}{9\left(1+\frac{x}{6}\right)^2}.
$$
Or, en dérivant le développement de $h(x)=\frac{1}{1+x}$, on obtient
$$
\frac{1}{(1+x)^2}=\sum_{n=0}^{+\infty}(-1)^n(n+1)x^n,
\qquad |x|<1.
$$
En remplaçant $x$ par $\frac{x}{6}$, il vient
$$
\frac{1}{\left(1+\frac{x}{6}\right)^2}
=
\sum_{n=0}^{+\infty}(-1)^n(n+1)\left(\frac{x}{6}\right)^n.
$$
Donc
$$
f(x)^2
=
\frac19\sum_{n=0}^{+\infty}(-1)^n(n+1)\left(\frac{x}{6}\right)^n
=
\sum_{n=0}^{+\infty}\frac{(-1)^n(n+1)}{9\cdot 6^n}x^n,
\qquad |x|<6.
$$
Ainsi
$$
\boxed{f(x)^2=\sum_{n=0}^{+\infty}\frac{(-1)^n(n+1)}{9\cdot 6^n}x^n.}
$$
Le rayon de convergence de ces deux séries est
$$
\boxed{R=6.}
$$
}

\end{enumerate}

}